\begin{abcliste}
\abc A moving charge has a magnetic field like a short piece of current: a positive charge like a current along its motion, a negative one like a current against it. Right hand, thumb along the motion of the proton (to the right): above the proton the fingers point \result{\text{out of the page}} ($\odot$).

\abc The force on a charge is $\vec F = q\,\vec v\times\vec B$. For a positive charge (right-hand rule: thumb $\vec v$ to the right, index finger $\vec B$ out of the page) the force would point down. The electron is negative, so the force on it points the other way: \result{\text{up, away from the proton}}.

\abc The electron moving to the right is like a current to the left. Its field at the proton, below it, points out of the page as well, and the force $e\,\vec v_1\times\vec B$ on the proton points down, away from the electron. The magnetic forces \result{\text{push the particles apart}}, just as two antiparallel currents repel each other.
\end{abcliste}
